% ===== PRÉAMBULE CANONIQUE — à coller VERBATIM en tête de CHAQUE .tex =====
\documentclass[11pt,a4paper]{article}
\usepackage[utf8]{inputenc}\usepackage[T1]{fontenc}\usepackage{lmodern}
\usepackage[french]{babel}
\usepackage{amsmath,amssymb,mathtools}\usepackage{icomma}
\usepackage{siunitx}\sisetup{locale=FR}  % \qty{}{}, \si{}, \ang{}
\usepackage{xcolor}\usepackage{colortbl}
\usepackage{tikz}\usepackage{pgfplots}\pgfplotsset{compat=1.18}
\usepackage[siunitx,europeanresistors]{circuitikz} % circuits (élec.)
\usepackage[most]{tcolorbox}\usepackage{enumitem}\usepackage{array,booktabs}
\usepackage[a4paper,margin=2.2cm]{geometry}
\usetikzlibrary{arrows.meta,calc,angles,quotes,positioning,patterns,%
  backgrounds,shapes.geometric,decorations.pathmorphing,decorations.markings,intersections}
\definecolor{primary}{RGB}{222,49,99}\definecolor{primarydark}{RGB}{150,22,55}
\definecolor{defcol}{RGB}{0,114,178}\definecolor{methcol}{RGB}{52,92,148}
\definecolor{excol}{RGB}{0,135,90}\definecolor{attncol}{RGB}{200,40,55}
\tcbset{boxbase/.style={enhanced,breakable,boxrule=0.9pt,arc=2.5pt,
  left=3mm,right=3mm,top=2.3mm,bottom=2.3mm,fonttitle=\bfseries,coltitle=white}}
% --- boîtes TUTORIEL (noms EXACTS, ne pas renommer) ---
\newtcolorbox{cle}[1][]{boxbase,colback=defcol!6,colframe=defcol,
  title={\textsc{À retenir}\ifx\relax#1\relax\relax\else\textmd{ : \itshape #1}\fi}}
\newtcolorbox{methode}[1][]{boxbase,colback=methcol!7,colframe=methcol,sharp corners,
  title={\textsc{Méthode}\ifx\relax#1\relax\relax\else\textmd{ --- \itshape #1}\fi}}
\newtcolorbox{exemplebox}[1][]{boxbase,colback=excol!5,colframe=excol,
  title={\textsc{Exemple}\ifx\relax#1\relax\relax\else\textmd{ : \itshape #1}\fi}}
\newtcolorbox{piege}[1][]{boxbase,colback=attncol!6,colframe=attncol,
  title={\textsc{Piège}\ifx\relax#1\relax\relax\else\textmd{ : \itshape #1}\fi}}
\newtcolorbox{remarque}[1][]{boxbase,colback=black!4,colframe=black!45,
  title={\textsc{Remarque}\ifx\relax#1\relax\relax\else\textmd{ : \itshape #1}\fi}}
% --- boîtes EXERCICE / CORRIGÉ (noms EXACTS, auto-comptées) ---
\newtcolorbox[auto counter]{exo}[1][]{boxbase,colback=primary!4,colframe=primary,
  title={\textsc{Exercice}~\thetcbcounter\ifx\relax#1\relax\relax\else\textmd{ --- \itshape #1}\fi}}
\newtcolorbox[auto counter]{corrige}[1][]{boxbase,colback=excol!4,colframe=excol,
  title={\textsc{Corrigé}~\thetcbcounter\ifx\relax#1\relax\relax\else\textmd{ --- \itshape #1}\fi}}
\newcommand{\vv}[1]{\vec{#1}}\newcommand{\ds}{\displaystyle}
\newcommand{\R}{\mathbb{R}}\newcommand{\N}{\mathbb{N}}\newcommand{\Z}{\mathbb{Z}}
% =========================================================================
\begin{document}
\usetikzlibrary{babel}
\begin{corrige}[combien d'électrons par seconde ?]
\begin{enumerate}[label=\alph*)]
  \item En $t=\qty{1}{\second}$, la charge est $Q=I\,t=3{,}2\times1=\qty{3.2}{\coulomb}$, soit
        \[ n=\frac{Q}{e}=\frac{3{,}2}{1{,}6\times10^{-19}}=2{,}0\times10^{19}\ \text{électrons}. \]
  \item Rapport à une mole : $\dfrac{2{,}0\times10^{19}}{6{,}02\times10^{23}}\approx3{,}3\times10^{-5}$.
        C'est donc \textbf{bien moins} qu'une mole d'électrons (environ un trente-millième de mole),
        alors même que le nombre d'électrons est gigantesque à notre échelle.
\end{enumerate}
\end{corrige}

\begin{corrige}[répartition dans trois branches en parallèle]
\begin{enumerate}[label=\alph*)]
  \item Loi des nœuds : $I=I_1+I_2+I_3$, donc
        \[ I_2=I-I_1-I_3=6-1{,}5-2=\qty{2.5}{\ampere}. \]
  \item L'ampèremètre principal mesure l'intensité \emph{totale} entrant au nœud $A$, soit
        $I=\qty{6}{\ampere}$.
\end{enumerate}
\end{corrige}

\begin{corrige}[un nœud, puis un fil coupé]
Au nœud $A$, le courant $I$ se partage entre la branche du haut ($I_1$) et celle du bas ($I_2$) :
$I=I_1+I_2$.
\begin{enumerate}[label=\alph*)]
  \item $I_1=I-I_2=2-0{,}8=\qty{1.2}{\ampere}$.
  \item Un fil coupé interrompt la branche du bas : plus aucun courant n'y circule, donc
        $I_2=\qty{0}{\ampere}$. Tout le courant passe alors par $L_1$ :
        $I_1=I-I_2=2-0=\qty{2}{\ampere}$.
\end{enumerate}
\end{corrige}

\begin{corrige}[intensité moyenne d'un courant variable]
\begin{enumerate}[label=\alph*)]
  \item $\ds i_{\text{moy}}=\frac{\Delta q}{\Delta t}=\frac{0{,}05}{0{,}2}=\qty{0.25}{\ampere}$.
  \item On fait tendre l'intervalle vers zéro : l'intensité instantanée est la dérivée de la charge
        par rapport au temps, $i(t)=\ds\lim_{\Delta t\to0}\frac{\Delta q}{\Delta t}=\frac{\mathrm{d}q}{\mathrm{d}t}$.
\end{enumerate}
\end{corrige}

\begin{corrige}[nœud avec courants entrants et sortants]
La loi des nœuds égale les intensités \emph{arrivantes} et \emph{partantes} :
\[ I_1+I_2=I_3+I_4 \;\Longrightarrow\; I_4=I_1+I_2-I_3=4+3-5=\qty{2}{\ampere}. \]
\end{corrige}

\end{document}
